% ============================================================================
% V. THỰC NGHIỆM
% ============================================================================
\section{THỰC NGHIỆM}

\subsection{Dữ liệu thực nghiệm từ hai bộ Parser mẫu}
Để kiểm định RAG framework sau khi lắp hai bộ \texttt{Parser} mẫu ở Mục IV.5, thực nghiệm được tiến hành trên toàn bộ hai kho tài liệu kỹ thuật chính thức: \textbf{Next.js Documentation} ở định dạng \texttt{.mdx}, \texttt{.md} và \textbf{Node.js Official API Documentation} ở định dạng \texttt{.md}. Thống kê quy mô và phân phối cấu trúc của 167 tài liệu gồm 162.757 dòng raw text, sinh ra 5.633 \texttt{chunk} sau khi chạy qua hai bộ \texttt{Parser} được trình bày tại Bảng \ref{tab:corpus_stats} và Bảng \ref{tab:chunk_dist}.

\begin{table}[H]
\centering
\caption{Thống kê kho ngữ liệu thực nghiệm sau khi indexing qua hai bộ \texttt{Parser} Next.js và Node.js}
\label{tab:corpus_stats}
\footnotesize\setlength{\tabcolsep}{3.5pt}\begin{tabularx}{\textwidth}{p{1.95cm}crrrrr>{\raggedright\arraybackslash}X}
\toprule
\textbf{Miền tài liệu} & \textbf{Định dạng} & \textbf{Tài liệu} & \textbf{Total Lines} & \textbf{Chunks} & \textbf{With Code} & \textbf{With Table} & \textbf{Scope \& Stability Distribution} \\
\midrule
Next.js Docs & \texttt{.mdx, .md} & 101 & 39.009 & 1.239 & 671 (54,16\%) & 106 (8,56\%) & Shared: 1.119; App Router: 95; Pages Router: 25 \\
\addlinespace[2pt]
Node.js API & \texttt{.md} & 66 & 123.748 & 4.394 & 1.436 (32,68\%) & 685 (15,59\%) & Stability 2: 2.974; Stability 1.x: 389; Stability 0: 161; Stability 3: 33; Khác: 837 \\
\midrule
\textbf{Tổng cộng} & -- & \textbf{167} & \textbf{162.757} & \textbf{5.633} & \textbf{2.107 (37,40\%)} & \textbf{791 (14,04\%)} & Chuẩn hóa \texttt{ltree} \& bảo toàn dòng 1:1 \\
\bottomrule
\end{tabularx}
\end{table}

\begin{table}[H]
\centering
\caption{Phân phối độ dài ký tự, số dòng và độ sâu cây phân cấp \texttt{ltree} của 5.633 chunk}
\label{tab:chunk_dist}
\footnotesize\setlength{\tabcolsep}{4pt}\begin{tabularx}{\textwidth}{>{\raggedright\arraybackslash}Xrrrrrrrc}
\toprule
\multirow{2}{*}{\textbf{Miền tài liệu}} & \multicolumn{5}{c}{\textbf{Độ dài $\text{text}^{\text{verb}}$ tính theo ký tự}} & \multicolumn{2}{c}{\textbf{Lines / Chunk}} & \textbf{Độ sâu \texttt{nlevel}} \\
\cmidrule(lr){2-6}\cmidrule(lr){7-8}\cmidrule(lr){9-9}
 & \textbf{Min} & \textbf{Median} & \textbf{Mean} & \textbf{P95} & \textbf{Max} & \textbf{Median} & \textbf{Mean / Max} & \textbf{Min / Med / Mean / Max} \\
\midrule
Next.js: 1.239 chunks & 60 & 890 & 1.185,6 & 3.112 & 8.787 & 18 & 29,5 / 325 & 3 / 5 / 4,80 / 7 \\
Node.js: 4.394 chunks & 40 & 356 & 729,7 & 2.518 & 15.726 & 16 & 26,3 / 637 & 2 / 4 / 4,08 / 6 \\
\bottomrule
\end{tabularx}
\end{table}

Số liệu tại Bảng \ref{tab:chunk_dist} phản ánh đặc trưng cấu trúc của chiến lược AST section chunking: median chunk length khá gọn với 890 ký tự trên Next.js và 356 ký tự trên Node.js, nhưng phân phối có dạng heavy-tailed với P95 tương ứng là 3.112 và 2.518 ký tự, và cực đại lên tới 15.726 ký tự ở các mục đặc tả table lớn hoặc ví dụ code đầy đủ.

Để đánh giá khách quan năng lực truy xuất của công cụ \texttt{hybrid\_search} mà không bị thiên lệch do người viết mã tự đặt câu hỏi nhìn trước cơ sở dữ liệu, bộ dữ liệu kiểm thử \texttt{data/eval/retrieval.jsonl} gồm 100 ca kiểm thử được xây dựng qua quy trình \textbf{3 pha cách ly mù (blind-isolated phases)} với 6 tác tử độc lập:
\begin{enumerate}
    \item \textbf{Pha 1 -- Sinh câu hỏi từ đề cương chủ đề (2 tác tử độc lập):} Tác tử chỉ được đọc danh mục tiêu đề của 167 tài liệu (\texttt{corpus\_topics\_outline.md}), hoàn toàn không được truy cập mã nguồn hay nội dung phân đoạn trong CSDL, để biên soạn câu hỏi cho hai lớp tác vụ:
    \begin{itemize}
        \item \textbf{Test set 1 -- \texttt{TS1-Lexical-API} (\texttt{e001}--\texttt{e050}):} 50 câu hỏi (25 Next.js, 25 Node.js) tra cứu chính xác tên hàm, chữ ký nạp chồng (overloaded signatures), thuộc tính cấu hình \nolinkurl{next.config.js}, mã lỗi hệ thống hoặc CLI flags.
        \item \textbf{Test set 2 -- \texttt{TS2-Conceptual} (\texttt{e051}--\texttt{e100}):} 50 câu hỏi (25 Next.js, 25 Node.js) diễn đạt tình huống kỹ thuật tự nhiên về kiến trúc, cơ chế vận hành và các chủ đề trùng tên tiêu đề con giữa các phạm vi phân cấp khác nhau (như \textit{Caching}, \textit{Revalidation}, \textit{Server Actions}, \textit{Streams}, \textit{Cluster}, \textit{TLS Session Resumption}).
    \end{itemize}
    \item \textbf{Pha 2 -- Lập lời gọi \texttt{hybrid\_search} đầu tiên (2 tác tử độc lập):} Đóng vai \texttt{Coding Agent} ở thời điểm gọi công cụ lần đầu, tác tử chỉ được đọc câu hỏi từ Pha 1 và bản đặc tả giao diện MCP của \texttt{hybrid\_search} (kèm danh mục 167 \texttt{doc\_slug}) để tự quyết định bộ tham số \texttt{query}, \texttt{pattern}, \texttt{is\_regex}, \texttt{case\_sensitive}, \texttt{doc\_slugs} và \texttt{path\_prefix}.
    \item \textbf{Pha 3 -- Đối chiếu CSDL và gán nhãn Ground-Truth (2 tác tử độc lập):} Giữ nguyên tuyệt đối câu hỏi và tham số truy vấn của Pha 1--2, tác tử truy vấn trực tiếp bảng \texttt{chunks} trong PostgreSQL để xác định đường dẫn phân cấp \texttt{ltree} chuẩn trả lời trực tiếp câu hỏi (\texttt{expected}). Tổng cộng có \textbf{100 đường dẫn \texttt{ltree} đích} được gán nhãn và xác thực tồn tại 100\% trong CSDL (Bảng \ref{tab:test_sets}).
\end{enumerate}

\begin{table}[H]
\centering
\caption{Cấu trúc các bộ dữ liệu kiểm thử truy xuất (\texttt{data/eval/retrieval.jsonl})}
\label{tab:test_sets}
\small\setlength{\tabcolsep}{4.5pt}\begin{tabularx}{\textwidth}{llccc>{\raggedright\arraybackslash}X}
\toprule
\textbf{Test Set} & \textbf{Phạm vi miền} & \textbf{Số câu} & \textbf{Tổng đích} & \textbf{Đích/câu} & \textbf{Mục tiêu kiểm định chính} \\
\midrule
\texttt{TS1-Lexical-API} & Next.js (25) + Node.js (25) & 50 & 50 & 1,00 & Độ chính xác khi tìm ký hiệu API, CLI flags, overloaded signatures và khớp mẫu regex/trigram \\
\addlinespace[2pt]
\texttt{TS2-Conceptual} & Next.js (25) + Node.js (25) & 50 & 50 & 1,00 & Khả năng nắm bắt ngữ nghĩa kiến trúc và phân biệt phạm vi phân cấp \texttt{ltree} nhờ context prefix \\
\midrule
\textbf{Tổng cộng} & \textbf{167 tài liệu (5.633 chunks)} & \textbf{100} & \textbf{100} & \textbf{1,00} & \textbf{Đánh giá toàn diện kênh đôi (Dual-Channel Venn) và nghiên cứu Ablation} \\
\bottomrule
\end{tabularx}
\end{table}

\subsection{Thiết lập thực nghiệm và Kiểm thử chức năng}

\subsubsection{Tham số cấu hình của framework và mô hình}
Các hyperparameters của framework cùng cơ sở lựa chọn tương ứng được tổng hợp tại Bảng \ref{tab:hyperparams}.

\begin{table}[htbp]
\centering
\caption{Hyperparameters, cấu hình và cơ sở lựa chọn của RAG framework}
\label{tab:hyperparams}
\footnotesize\setlength{\tabcolsep}{4pt}\begin{tabularx}{\textwidth}{>{\raggedright\arraybackslash}p{3.2cm}>{\raggedright\arraybackslash}p{4.6cm}>{\raggedright\arraybackslash}X}
\toprule
\textbf{Tham số / Thành phần} & \textbf{Giá trị thiết lập} & \textbf{Cơ sở lựa chọn} \\
\midrule
Mô hình Dense Embedding & \texttt{Qwen/Qwen3-Embedding-0.6B} [4] & Hỗ trợ đa ngôn ngữ và tài liệu lập trình, vừa bộ nhớ GPU 4 GB ở chế độ FP16 \\
\addlinespace[2pt]
Số chiều Vector $d$ & 512 chiều, MRL truncation [5] + chuẩn hóa $L_2$ & Giảm 50\% dung lượng lưu trữ so với 1.024 chiều gốc mà vẫn giữ tương đồng Cosine \\
\addlinespace[2pt]
Vector Index \texttt{pgvector} & HNSW \texttt{vector\_cosine\_ops} (\texttt{m=16, ef\_construction=64}) & Cấu hình mặc định chuẩn của \texttt{pgvector} cân bằng tốc độ dựng chỉ mục và độ phủ trên 5.633 chunks \\
\addlinespace[2pt]
Chuỗi tối đa Embedding & $L_{\text{emb}} = 2.048$ tokens, \texttt{left-padding}, batch size = 8--16 (fallback 1) & Bao phủ trọn vẹn P95 độ dài chunk (3.112 ký tự) mà không tràn 4 GB VRAM \\
\addlinespace[2pt]
Cross-Encoder Reranker & \texttt{mmarco-mMiniLMv2-L12-H384} [8], $L_{\max}=256$, $\alpha=1,0$, $P=10$ & Mô hình 12 lớp gọn nhẹ giữ độ trễ rerank $< 400$ ms; LRU cache 2.048 cặp giảm tính toán lặp \\
\addlinespace[2pt]
Hybrid Search \& RRF & $k_{\text{rrf}} = 60$, bể ứng viên $C = \max(6K, 120)$, FTS \texttt{simple} & Hằng số $k_{\text{rrf}}=60$ theo Cormack et al. [6]; từ điển \texttt{simple} bảo toàn nguyên vẹn định danh API \\
\addlinespace[2pt]
Trọng số Sparse FTS & Proximity boost = $4,0$, exact match bonus = $+2,0$ & Ưu tiên các phân đoạn chứa cụm từ khóa đứng liền kề và hội đủ toàn bộ token truy vấn \\
\addlinespace[2pt]
Ngưỡng confidence $\gamma$ & $\tau_{\text{low}} = 0,70$, $\tau_{\text{high}} = 0,82$, $\tau_{\text{veto}} = -1,0$ & Phân tầng độ tin cậy thực nghiệm để kích hoạt cảnh báo khi thiếu bằng chứng khớp từ vựng \\
\addlinespace[2pt]
Giới hạn MCP \& Parser & Nav $\le 50$ nodes; Graph $\le 100$ bước; Raw $\le 200$ dòng; 5 invariants & Chặn tràn cửa sổ ngữ cảnh của \texttt{Coding Agent} và bảo toàn 100\% ranh giới cú pháp AST \\
\bottomrule
\end{tabularx}
\end{table}

\subsubsection{Cấu hình phần cứng và môi trường}
Môi trường phần cứng và phần mềm phục vụ thực nghiệm được tổng hợp tại Bảng \ref{tab:env}.

\begin{table}[htbp]
\centering
\caption{Cấu hình phần cứng và môi trường thực nghiệm}
\label{tab:env}
\small\setlength{\tabcolsep}{5pt}\begin{tabularx}{\textwidth}{p{4.6cm}>{\raggedright\arraybackslash}X}
\toprule
\textbf{Thành phần} & \textbf{Thông số kỹ thuật} \\
\midrule
CPU & AMD Ryzen 5 5500U with Radeon Graphics (6 nhân / 12 luồng xử lý) \\
\addlinespace[2pt]
GPU & NVIDIA GeForce GTX 1650 (4 GB VRAM, suy luận FP16 trên CUDA) \\
\addlinespace[2pt]
RAM \& Hệ điều hành & 22 GiB DDR4; Linux Ubuntu / Debian (x86\_64) \\
\addlinespace[2pt]
Hệ quản trị CSDL & PostgreSQL 16 kèm tiện ích mở rộng \texttt{pgvector}, \texttt{ltree}, \texttt{pg\_trgm} \\
\addlinespace[2pt]
Ngôn ngữ \& Thư viện chính & Python 3.13.11 quản lý qua \texttt{uv}, PyTorch, Sentence-Transformers, \texttt{asyncpg}, \texttt{markdown-it-py}, FastAPI \\
\addlinespace[2pt]
Thời gian nạp \& lập chỉ mục & Phân tích AST \& kiểm định 5 bất biến (167 tài liệu, 5.633 chunks): $\approx 2,1$ giây; Nhúng lạnh toàn kho (5.633 chunks trên GTX 1650 4 GB): $9.793,4$ giây ($\approx 1,74$ giây/chunk); Nạp ấm qua differential vector cache: $< 3,0$ giây \\
\addlinespace[2pt]
Độ trễ truy vấn (\texttt{p50} / \texttt{p95}) & Sparse FTS: $2,8\text{ ms} / 5,6\text{ ms}$; Dense Vector: $205,3\text{ ms} / 222,8\text{ ms}$; Hybrid RRF: $206,5\text{ ms} / 223,9\text{ ms}$; Dual-Channel Venn (không Rerank): $207,6\text{ ms} / 218,2\text{ ms}$; Full Dual-Channel Venn (có Rerank): $585,5\text{ ms} / 634,4\text{ ms}$ \\
\addlinespace[2pt]
Độ trễ khởi tạo phiên MCP & Subprocess \texttt{stdio} lạnh (nạp lại Qwen3 + Reranker): $\approx 29.727\text{ ms}$ ở lệnh đầu; Warm Streamable HTTP (\texttt{POST /api/mcp}) dùng chung GPU Engine: $\approx 695\text{--}768\text{ ms}$ ở lệnh đầu (giảm $\approx 40\times$) \\
\bottomrule
\end{tabularx}
\end{table}

\subsubsection{Kiểm thử chức năng và bất biến toàn vẹn dữ liệu}
Trước khi đánh giá chất lượng xếp hạng, hệ thống được kiểm tra tự động qua bộ test tích hợp (\texttt{uv run scripts/check.py}) và kiểm định trực tiếp trên toàn bộ 5.633 phân đoạn trong CSDL nhằm xác nhận các bất biến kiến trúc (Bảng \ref{tab:func_tests}).

\begin{table}[htbp]
\centering
\caption{Kết quả kiểm thử chức năng và kiểm định bất biến trên kho dữ liệu thực tế}
\label{tab:func_tests}
\footnotesize\setlength{\tabcolsep}{4pt}\begin{tabularx}{\textwidth}{c>{\raggedright\arraybackslash}p{4.3cm}>{\raggedright\arraybackslash}Xc}
\toprule
\textbf{Mã KT} & \textbf{Hạng mục kiểm thử} & \textbf{Phương pháp và phạm vi kiểm chứng} & \textbf{Kết quả} \\
\midrule
\texttt{KT01} & Bất biến ánh xạ dòng 1:1 và bảo toàn code block & Đối chiếu \texttt{start\_line}, \texttt{end\_line} của toàn bộ 5.633 chunks với \texttt{raw\_text} gốc trên 167 tài liệu; kiểm tra cân bằng dấu rào code trên 2.107 chunks chứa code. & Đạt (0 lỗi) \\
\addlinespace[2pt]
\texttt{KT02} & Bất biến đường dẫn \texttt{ltree} và tiền tố ngữ cảnh & Kiểm tra biểu thức chính quy nhãn \texttt{ltree} và tiền tố \texttt{[Tài liệu > Mục cha]} đầu trường \texttt{contextualized\_text} trên toàn bộ 5.633 chunks. & Đạt (100\%) \\
\addlinespace[2pt]
\texttt{KT03} & Cơ chế từ chối định dạng lạ (\texttt{CodingParserRegistry}) & Thử nạp tệp nằm ngoài đặc tả \texttt{NextJSMDXParser} và \texttt{NodeJSMDParser}; xác nhận hệ thống ném ngoại lệ \texttt{UnsupportedFrameworkError} tức thì, không dùng generic fallback. & Đạt \\
\addlinespace[2pt]
\texttt{KT04} & Nạp ấm bảo toàn bộ nhớ đệm (\texttt{Differential Vector Cache}) & Chạy lại quy trình nạp khi \texttt{contextualized\_text} không đổi; xác nhận hàm băm MD5 giữ nguyên 5.633 vector 512 chiều mà không gọi lại GPU. & Đạt ($< 3$s) \\
\addlinespace[2pt]
\texttt{KT05} & Kiểm chứng trích dẫn nguyên văn Knowledge Graph & Gọi \texttt{graph\_curate\_edges} với trích dẫn hợp lệ (mở rộng container node xuống các leaf nodes, ghi \texttt{graph\_edge\_log}) và trích dẫn bịa đặt (bị chặn bởi \texttt{E\_CITATION\_NOT\_FOUND}). & Đạt \\
\addlinespace[2pt]
\texttt{KT06} & Duyệt phân cấp \texttt{ltree} và trích xuất dòng thô & Kiểm tra \texttt{hierarchical\_navigate} (tái tạo virtual parent nodes cho mục container) và \texttt{get\_raw} (chặn biên nghiêm ngặt tối đa 200 dòng/lần). & Đạt \\
\addlinespace[2pt]
\texttt{KT07} & Kiểm chứng trích dẫn câu trả lời (\texttt{submit\_grounded\_answer}) & Xác thực mọi trích dẫn \texttt{(path, start\_line, end\_line)} của Agent trên PostgreSQL trước khi chấp nhận câu trả lời; từ chối ngay khi đường dẫn \texttt{ltree} không tồn tại hoặc thiếu trích dẫn. & Đạt \\
\bottomrule
\end{tabularx}
\end{table}

\subsection{Evaluation Metrics và Quy tắc chấm điểm}
Trên bộ kiểm thử gồm $|Q| = 100$ truy vấn, mỗi truy vấn $q \in Q$ có tập nhãn đích $G_q \subset \mathcal{C}$ gồm đường dẫn phân cấp \texttt{ltree} chuẩn. Do tài liệu kỹ thuật có cấu trúc cây phân cấp sâu (tới 7 tầng ở Next.js và 6 tầng ở Node.js), một phân đoạn trả về có đường dẫn $p$ được tính là trả lời đúng cho nhãn đích $e \in G_q$ nếu $p$ trùng khớp hoàn toàn hoặc là nút con trực tiếp nằm trong nhánh của $e$:
\begin{equation}
    \text{answers}(p, G_q) \iff \exists\, e \in G_q : (p = e) \lor \text{startswith}(p,\, e + \text{``\texttt{.}''})
    \label{eq:ltree_match}
\end{equation}

Gọi $\hat{R}_q^{(k)} = (\hat{c}_1, \hat{c}_2, \dots, \hat{c}_k)$ là danh sách $k$ phân đoạn đầu tiên hệ thống trả về (sắp theo thứ tự ưu tiên giỏ Venn: \texttt{both} $\to$ \texttt{semantic\_only} $\to$ \texttt{verbatim\_only}):
\begin{itemize}
    \item \textbf{Tỷ lệ tìm thấy (\texttt{Found}) và phân rã giỏ Venn:} Một truy vấn được tính là \texttt{Found} nếu tồn tại ít nhất một phân đoạn thỏa mãn Phương trình (\ref{eq:ltree_match}) trong kết quả hợp nhất của \texttt{hybrid\_search}, đồng thời được phân vào một trong bốn giỏ loại trừ nhau: \textbf{Cả hai kênh (\texttt{both})}, \textbf{Chỉ ngữ nghĩa (\texttt{semantic\_only})}, \textbf{Chỉ nguyên văn (\texttt{verbatim\_only})} hoặc \textbf{Bỏ sót (\texttt{missed})}.
    \item \textbf{Hit@k và Recall@k:}
    \begin{equation}
    \begin{aligned}
        \text{Hit@}k &= \frac{1}{|Q|} \sum_{q \in Q} \mathbb{I}\!\left(\exists\, \hat{c} \in \hat{R}_q^{(k)} : \text{answers}(\hat{c}, G_q)\right), \\
        \text{Recall@}k &= \frac{1}{|Q|} \sum_{q \in Q} \frac{\left|\left\{e \in G_q \mid \exists\, \hat{c} \in \hat{R}_q^{(k)} : \text{answers}(\hat{c}, \{e\})\right\}\right|}{|G_q|}
    \end{aligned}
    \label{eq:recall_at_k}
    \end{equation}
    \item \textbf{MRR@k -- Mean Reciprocal Rank:} Với $r_q^*$ là thứ hạng đầu tiên xuất hiện một phân đoạn thỏa mãn $\text{answers}(\hat{c}_i, G_q)$:
    \begin{equation}
        \text{MRR@}k = \frac{1}{|Q|} \sum_{q \in Q} \frac{1}{r_q^*}, \quad r_q^* = \min\left\{i \in \{1, \dots, k\} \mid \text{answers}(\hat{c}_i, G_q)\right\}, \quad \frac{1}{\infty} = 0
        \label{eq:mrr}
    \end{equation}
    \item \textbf{NDCG@k -- Normalized Discounted Cumulative Gain:}
    \begin{equation}
        \text{NDCG@}k = \frac{1}{|Q|} \sum_{q \in Q} \frac{\text{DCG}_q\text{@}k}{\text{IDCG}_q\text{@}k}, \qquad
        \text{DCG}_q\text{@}k = \sum_{i=1}^{k} \frac{\mathbb{I}_{\text{unique}}(\hat{c}_i, G_q)}{\log_2(i + 1)}
        \label{eq:ndcg}
    \end{equation}
    \item \textbf{Kiểm định McNemar hai phía chính xác (Exact Two-Sided McNemar Test):} Để kiểm chứng sự khác biệt giữa hai cấu hình $A$ và $B$ trên cùng $100$ truy vấn có phải do ngẫu nhiên hay không, gọi $b$ là số câu $A$ đúng nhưng $B$ sai (\textit{Tốt lên}) và $c$ là số câu $A$ sai nhưng $B$ đúng (\textit{Kém đi}). Giá trị $p\text{-value}$ nhị thức hai phía chính xác được tính theo:
    \begin{equation}
        p\text{-value} = \min\!\left(1,\; 2 \sum_{i=0}^{\min(b, c)} \binom{b+c}{i} 0,5^{b+c}\right)
        \label{eq:mcnemar}
    \end{equation}
\end{itemize}

\subsection{Kết quả thực nghiệm trên từng test set}

\subsubsection{Kết quả chi tiết theo test set, miền tài liệu và kênh Venn}
Bảng \ref{tab:results_per_testset} ghi nhận kết quả của cấu hình \textbf{Dual-Channel Venn} trên hai bộ test \texttt{TS1-Lexical-API} và \texttt{TS2-Conceptual}, tách riêng theo từng miền tài liệu \texttt{Next.js} và \texttt{Node.js}.

\begin{table}[htbp]
\centering
\caption{Kết quả truy xuất của cấu hình Dual-Channel Venn theo từng bộ test và miền dữ liệu (\%)}
\label{tab:results_per_testset}
\scriptsize\setlength{\tabcolsep}{3pt}\begin{tabularx}{\textwidth}{>{\raggedright\arraybackslash}Xlcccccccc}
\toprule
\textbf{Test Set} & \textbf{Miền} & \textbf{Found} & \textbf{Giỏ (\texttt{B}/\texttt{S}/\texttt{V}/\texttt{M})} & \textbf{Hit@1} & \textbf{Hit@5} & \textbf{R@5} & \textbf{R@10} & \textbf{MRR@10} & \textbf{NDCG@10} \\
\midrule
\multirow{3}{*}{\texttt{TS1-Lexical-API}} & Next.js (25) & 96,0 & 20 / 3 / 1 / 1 & 72,0 & 88,0 & 90,0 & 92,0 & 78,2 & 84,9 \\
 & Node.js (25) & \textbf{100,0} & 24 / 1 / 0 / 0 & 76,0 & \textbf{100,0} & \textbf{100,0} & \textbf{100,0} & 87,3 & 93,6 \\
\cmidrule(lr){2-10}
 & \textbf{Tổng TS1 (50)} & \textbf{98,0} & \textbf{44 / 4 / 1 / 1} & \textbf{74,0} & \textbf{94,0} & \textbf{95,0} & \textbf{96,0} & \textbf{82,8} & \textbf{89,2} \\
\midrule
\multirow{3}{*}{\texttt{TS2-Conceptual}} & Next.js (25) & 92,0 & 19 / 4 / 0 / 2 & 44,0 & 88,0 & 88,0 & 92,0 & 63,0 & 71,8 \\
 & Node.js (25) & \textbf{100,0} & 22 / 2 / 1 / 0 & 64,0 & 92,0 & 92,0 & 96,0 & 77,4 & 83,0 \\
\cmidrule(lr){2-10}
 & \textbf{Tổng TS2 (50)} & \textbf{96,0} & \textbf{41 / 6 / 1 / 2} & \textbf{54,0} & \textbf{90,0} & \textbf{90,0} & \textbf{94,0} & \textbf{70,2} & \textbf{77,4} \\
\midrule
\multicolumn{2}{l}{\textbf{Toàn cục (100 câu, có Rerank)}} & \textbf{97,0} & \textbf{85 / 10 / 2 / 3} & \textbf{64,0} & \textbf{92,0} & \textbf{92,5} & \textbf{95,0} & \textbf{76,5} & \textbf{83,3} \\
\multicolumn{2}{l}{\textbf{Toàn cục (100 câu, không Rerank)}} & \textbf{97,0} & \textbf{85 / 10 / 2 / 3} & \textbf{68,0} & \textbf{92,0} & \textbf{92,5} & \textbf{95,0} & \textbf{78,9} & \textbf{85,0} \\
\bottomrule
\end{tabularx}
\end{table}

Số liệu tại Bảng \ref{tab:results_per_testset} cho thấy ba đặc trưng quan trọng:
\begin{enumerate}
    \item \textbf{Độ phủ 97,0\% trên toàn tập và 100,0\% khi phạm vi lọc hợp lệ:} Cả hai biến thể của Dual-Channel Venn đều tìm thấy phân đoạn đích trên \textbf{97/100 câu hỏi (Found = 97,0\%)}. Đáng chú ý, trên miền \texttt{Node.js} (50 câu), tỷ lệ tìm thấy đạt tuyệt đối \textbf{100,0\%} ở cả hai bộ \texttt{TS1} và \texttt{TS2}. Ba câu duy nhất bị bỏ sót trên \texttt{Next.js} (\texttt{e011}, \texttt{e070}, \texttt{e074}) đều do tác tử Pha 2 truyền bộ lọc cứng \texttt{doc\_slugs} chỉ gồm 3 tài liệu ứng viên không chứa tài liệu đích; trên toàn bộ \textbf{97/97 câu hỏi mà \texttt{doc\_slugs} bao hàm tài liệu đích, hệ thống đạt tỷ lệ tìm thấy 100\%}, trong đó \textbf{85 câu} khớp đồng thời cả hai kênh (\texttt{both}), \textbf{10 câu} tìm thấy qua kênh ngữ nghĩa (\texttt{semantic\_only}), và \textbf{2 câu} (\texttt{e021}, \texttt{e083}) được cứu sống hoàn toàn nhờ kênh nguyên văn (\texttt{verbatim\_only}).
    \item \textbf{Độ chính xác rất cao trên các truy vấn tra cứu API (\texttt{TS1-Lexical-API}):} Với các câu hỏi nhắm vào tên hàm, chữ ký nạp chồng và cờ cấu hình, hệ thống đạt \textbf{Recall@5 = 95,0\%}, \textbf{Recall@10 = 96,0\%}, \textbf{Hit@1 = 74,0\%}, \textbf{MRR@10 = 82,8\%--83,4\%} và \textbf{NDCG@10 = 89,2\%--89,9\%}. Riêng trên miền \texttt{Node.js} của \texttt{TS1}, hệ thống đạt tuyệt đối \textbf{Found = 100,0\%, Hit@5 = 100,0\%, Recall@5 = 100,0\% và Recall@10 = 100,0\%}.
    \item \textbf{Khả năng bao phủ các truy vấn tình huống kiến trúc (\texttt{TS2-Conceptual}):} Dù câu hỏi trong \texttt{TS2} mô tả hiện tượng kỹ thuật dài dòng mà không nêu trực tiếp tiêu đề mục, hệ thống vẫn đạt \textbf{Recall@5 = 90,0\%} và \textbf{Recall@10 = 94,0\%} (trên \texttt{Node.js} đạt \textbf{Recall@10 = 96,0\%}).
\end{enumerate}

\subsubsection{Phân bố nhãn độ tin cậy (\texorpdfstring{$\gamma$}{gamma}) của cổng Sparse Match Gating}
Bảng \ref{tab:confidence_eval} đối chiếu phân bố nhãn độ tin cậy $\gamma \in \{\texttt{none}, \texttt{low}, \texttt{medium}, \texttt{high}\}$ giữa các cấu hình truy xuất đơn lẻ và cấu hình kết hợp.

\begin{table}[htbp]
\centering
\caption{Phân bố nhãn độ tin cậy $\gamma$ qua cổng kiểm soát Sparse Match Gating trên 100 truy vấn}
\label{tab:confidence_eval}
\footnotesize\setlength{\tabcolsep}{4pt}\begin{tabularx}{\textwidth}{lcccc>{\raggedright\arraybackslash}X}
\toprule
\textbf{Cấu hình / Tập truy vấn} & \textbf{\texttt{none} (\%)} & \textbf{\texttt{low} (\%)} & \textbf{\texttt{medium} (\%)} & \textbf{\texttt{high} (\%)} & \textbf{Nhận xét hành vi cổng kiểm soát} \\
\midrule
1. Chỉ Sparse FTS (\texttt{ts\_rank}) & 3,0 & 97,0 & 0,0 & 0,0 & Thiếu tương đồng Cosine ($s_{\max}=0$) nên kẹp ở mức \texttt{low} \\
\addlinespace[2pt]
2. Chỉ Dense Vector (\texttt{Qwen3-0.6B}) & 24,0 & 76,0 & 0,0 & 0,0 & Thiếu khớp từ vựng ($\text{has\_sparse}=\text{False}$) nên kẹp ở \texttt{low} \\
\addlinespace[2pt]
6. Full Dual-Channel Venn (\texttt{TS1-API}) & 0,0 & 30,0 & 56,0 & 14,0 & 70,0\% truy vấn API đạt mức tin cậy \texttt{medium} hoặc \texttt{high} \\
\addlinespace[2pt]
6. Full Dual-Channel Venn (\texttt{TS2-Concept}) & 0,0 & 26,0 & 66,0 & 8,0 & 74,0\% truy vấn khái niệm đạt mức \texttt{medium} hoặc \texttt{high} \\
\midrule
\textbf{6. Toàn cục (100 truy vấn)} & \textbf{0,0} & \textbf{28,0} & \textbf{61,0} & \textbf{11,0} & \textbf{72,0\%--73,0\% đạt mức \texttt{medium}/\texttt{high}; 0\% bị gán \texttt{none}} \\
\bottomrule
\end{tabularx}
\end{table}

\subsection{Đối sánh các cấu hình và Ablation Study}
Để tách biệt đóng góp thực tế của từng mô-đun trong framework, nghiên cứu Ablation được thực hiện trên cùng 100 ca kiểm thử qua 6 cấu hình tại Bảng \ref{tab:comparison}:
\begin{enumerate}
    \item \textbf{AST + Sparse FTS Only}: Chỉ xếp hạng bằng điểm từ vựng $s_{\text{sparse}}(q, c)$ ở Phương trình (\ref{eq:sparse}) với cấu hình \texttt{simple}.
    \item \textbf{AST + Dense Vector Only}: Chỉ xếp hạng bằng độ tương đồng Cosine $s_{\text{dense}}(q, c)$ ở Phương trình (\ref{eq:cosine}) trên vector 512 chiều \texttt{Qwen3-Embedding-0.6B}.
    \item \textbf{AST + Hybrid RRF w/o Rerank}: Kết hợp Dense + Sparse trên $\text{text}^{\text{ctx}}_c$ (đã gắn tiền tố ngữ cảnh cha \texttt{[Tài liệu > Mục cha]}) qua công thức RRF ở Phương trình (\ref{eq:rrf}) với $k_{\text{rrf}}=60$, không dùng Cross-Encoder.
    \item \textbf{AST + Hybrid RRF + Cross-Encoder Rerank}: Bật thêm tầng xếp hạng lại bằng \texttt{mmarco-mMiniLMv2} ($L_{\max} = 256$ token, $\alpha = 1,0$) trên $P = 10$ ứng viên đầu tiên của kênh ngữ nghĩa.
    \item \textbf{Dual-Channel Venn w/o Rerank}: Hợp nhất theo phân hoạch Venn giữa kênh ngữ nghĩa Hybrid RRF (không Rerank) và kênh nguyên văn regex/trigram.
    \item \textbf{Full Dual-Channel Venn (with Rerank)}: Cấu hình đầy đủ kết hợp kênh ngữ nghĩa đã qua Cross-Encoder Rerank với kênh nguyên văn regex/trigram và gom nhóm theo ngữ cảnh cha.
\end{enumerate}

\begin{table}[htbp]
\centering
\caption{Đối sánh các cấu hình Ablation và độ trễ truy vấn trên 100 ca kiểm thử (\texttt{data/eval/ablation\_results.json})}
\label{tab:comparison}
\scriptsize\setlength{\tabcolsep}{3pt}\begin{tabularx}{\textwidth}{>{\raggedright\arraybackslash}Xccccccccc}
\toprule
\multirow{2}{*}{\textbf{Phương pháp / Cấu hình}} & \multirow{2}{*}{\textbf{Found}} & \multicolumn{2}{c}{\textbf{MRR@10 theo bộ (\%)}} & \multicolumn{5}{c}{\textbf{Toàn cục trên 100 truy vấn (\%)}} & \textbf{Latency (ms)} \\
\cmidrule(lr){3-4}\cmidrule(lr){5-9}\cmidrule(lr){10-10}
 & \textbf{(\%)} & \textbf{\texttt{TS1-API}} & \textbf{\texttt{TS2-Con}} & \textbf{Hit@1} & \textbf{Hit@5} & \textbf{R@5} & \textbf{NDCG@10} & \textbf{MRR@10} & \textbf{p50 / p95} \\
\midrule
1. AST + Sparse FTS Only & 88,0 & 77,8 & 67,9 & 60,0 & 86,0 & 86,0 & 78,6 & 72,9 & 2,8 / 5,6 \\
2. AST + Dense Vector Only & 95,0 & 70,6 & 73,1 & 59,0 & 90,0 & 91,0 & 79,6 & 71,9 & 205,3 / 222,8 \\
3. AST + Hybrid RRF w/o Rerank & 95,0 & 82,8 & \textbf{75,2} & \textbf{68,0} & \textbf{93,0} & \textbf{93,0} & \textbf{85,0} & \textbf{79,0} & 206,5 / 223,9 \\
4. AST + Hybrid RRF + Rerank & 95,0 & 74,6 & 66,9 & 58,0 & 89,0 & 89,5 & 79,0 & 70,7 & 584,6 / 619,9 \\
5. Dual-Channel Venn w/o Rerank & \textbf{97,0} & \textbf{83,4} & 74,4 & \textbf{68,0} & 92,0 & 92,5 & \textbf{85,0} & 78,9 & 207,6 / 218,2 \\
\textbf{6. Full Dual-Channel Venn} & \textbf{97,0} & 82,8 & 70,2 & 64,0 & 92,0 & 92,5 & 83,3 & 76,5 & 585,5 / 634,4 \\
\bottomrule
\end{tabularx}
\end{table}

Để kiểm chứng ý nghĩa thống kê của các mức chênh lệch giữa từng cặp cấu hình trong Bảng \ref{tab:comparison}, Bảng \ref{tab:mcnemar} trình bày kết quả kiểm định McNemar hai phía chính xác theo Phương trình (\ref{eq:mcnemar}).

\begin{table}[htbp]
\centering
\caption{Kiểm định McNemar hai phía chính xác giữa các cặp cấu hình Ablation ($N = 100$)}
\label{tab:mcnemar}
\footnotesize\setlength{\tabcolsep}{3.2pt}\begin{tabularx}{\textwidth}{>{\raggedright\arraybackslash}Xccccl}
\toprule
\textbf{Cặp đối sánh (Cấu hình $A$ vs. $B$)} & \textbf{Chỉ số} & \textbf{Tốt lên ($b$)} & \textbf{Kém đi ($c$)} & \textbf{$p$-value} & \textbf{Kết luận thống kê ($\alpha = 0,05$)} \\
\midrule
Hybrid RRF (3) vs. Sparse FTS Only (1) & \texttt{Found / Hit@5} & 7 & 0 & \textbf{0,0156} & Có ý nghĩa thống kê ($p < 0,05$) \\
Hybrid RRF (3) vs. Sparse FTS Only (1) & \texttt{Hit@1} & 12 & 4 & 0,0768 & Cải thiện rõ rệt ở biên $\alpha = 0,10$ \\
Hybrid RRF (3) vs. Dense Vector Only (2) & \texttt{Hit@1} & 17 & 8 & 0,1078 & Nâng \texttt{Hit@1} từ $59,0\%$ lên $68,0\%$ (+9,0 điểm \%) \\
Sparse FTS Only (1) vs. Dense Vector Only (2) & \texttt{Found} & 2 & 9 & 0,0654 & Bổ trợ ngược chiều giữa \texttt{TS1} và \texttt{TS2} \\
Dual-Channel (6) vs. Hybrid + Rerank (4) & \texttt{Hit@1} & 10 & 4 & 0,1796 & Kéo 10 câu bị Reranker hạ bậc về lại hạng 1 \\
Dual-Channel (6) vs. Hybrid RRF (3, 4) & \texttt{Found} & 2 & 0 & 0,5000 & Cứu 2 câu \texttt{e021}, \texttt{e083} nhờ giỏ \texttt{verbatim\_only} \\
\bottomrule
\end{tabularx}
\end{table}

\subsection{Thảo luận, đánh đổi kiến trúc và phân tích lỗi chi tiết}
Từ số liệu thực nghiệm tại Bảng \ref{tab:results_per_testset}--\ref{tab:mcnemar} và kết quả đối chiếu từng ca truy vấn cụ thể trong tệp \nolinkurl{ablation_results.json}, có thể rút ra năm nhận định kỹ thuật cốt lõi:

\begin{itemize}
    \item \textbf{Sự bổ trợ ngược chiều giữa Sparse FTS và Dense Vector:} Khi đứng riêng lẻ, \textbf{Sparse FTS (1)} đạt \texttt{MRR@10 = 77,8\%} trên \texttt{TS1-Lexical-API} nhờ khớp chính xác tên định danh, nhưng tụt xuống \texttt{67,9\%} trên \texttt{TS2-Conceptual} và bỏ sót tới 12/100 câu (\texttt{Found = 88,0\%}) khi câu hỏi mô tả hiện tượng bằng từ vựng khác với tài liệu hoặc khi từ điển \texttt{simple} không tách các biểu thức có dấu câu đặc thù. Ngược lại, \textbf{Dense Vector (2)} nâng tỷ lệ tìm thấy lên \texttt{95,0\%} (\texttt{Recall@10 = 95,5\%}) và đạt \texttt{MRR@10 = 73,1\%} trên \texttt{TS2-Conceptual}, nhưng chỉ đạt \texttt{MRR@10 = 70,6\%} trên \texttt{TS1}. Khi hợp nhất qua \textbf{Hybrid RRF (3)}, \texttt{Hit@1} tăng từ $60,0\%$ (Sparse) và $59,0\%$ (Dense) lên \textbf{68,0\%}, \texttt{Recall@5} đạt \textbf{93,0\%}, \texttt{MRR@10} tăng lên \textbf{79,0\%} và \texttt{NDCG@10} đạt \textbf{85,0\%}, vượt trội có ý nghĩa thống kê so với Sparse FTS trên cả \texttt{Found} và \texttt{Hit@5} ($7$ câu tốt lên so với $0$ câu kém đi, $p = 0,0156$).
    \item \textbf{Sự đánh đổi thực tế giữa bảo toàn atomic AST chunk và cửa sổ 256 token của Cross-Encoder:} Bảng \ref{tab:comparison} phơi bày một hiện tượng thực nghiệm rất đáng chú ý khi đối chiếu cấu hình \textbf{(3) Hybrid RRF w/o Rerank} và \textbf{(4) Hybrid RRF + Cross-Encoder Rerank}:
    \begin{enumerate}
        \item Tầng Cross-Encoder đẩy thành công \textbf{10 câu} từ các vị trí bên dưới lên thẳng hạng 1 (\texttt{e003}, \texttt{e024}, \texttt{e036}, \texttt{e038}, \texttt{e047}, \texttt{e051}, \texttt{e060}, \texttt{e061}, \texttt{e085}, \texttt{e089}) khi thông tin trả lời nằm gọn trong 256 token đầu tiên của phân đoạn.
        \item Tuy nhiên, nó lại làm tụt hạng 1 của \textbf{20 câu khác} (khiến \texttt{Hit@1} của riêng nhánh ngữ nghĩa giảm từ $68,0\%$ xuống $58,0\%$, \texttt{MRR@10} giảm từ $79,0\%$ xuống $70,7\%$), đồng thời làm độ trễ trung vị \texttt{p50} tăng từ $206,5\text{ ms}$ lên $584,6\text{ ms}$. Phân tích các câu bị tụt hạng sâu như câu \texttt{e078} (\nolinkurl{nodejs.child_process.overview}: tụt từ hạng 1 xuống hạng 9), câu \texttt{e068} (tụt từ hạng 1 xuống hạng 8), câu \texttt{e029} (tụt từ hạng 1 xuống hạng 7), câu \texttt{e012}, \texttt{e054}, \texttt{e076} (tụt từ hạng 1 xuống hạng 5) hay câu \texttt{e086}, \texttt{e098} (tụt từ hạng 1 xuống hạng 4) chỉ ra nguyên nhân trực tiếp: do chiến lược AST chunking giữ nguyên vẹn toàn bộ mục tài liệu thành một phân đoạn đơn lẻ (P95 dài tới $2.518$--$3.112$ ký tự), phần giải thích chi tiết trả lời cho câu hỏi thường nằm ở nửa sau của phân đoạn, vượt ra ngoài cửa sổ cắt $L_{\max} = 256$ token của \texttt{mmarco-mMiniLMv2}. Khi sử dụng hệ số trộn $\alpha = 1,0$, điểm Cross-Encoder tính trên 256 token đầu đã ghi đè lên thứ hạng RRF vốn được tính trên toàn bộ 2.048 token.
    \end{enumerate}
    \item \textbf{Cơ chế tự sửa lỗi và cứu hộ của kênh nguyên văn (Dual-Channel Venn):} Việc hợp nhất song song kênh nguyên văn regex/trigram ở cấu hình \textbf{(5, 6)} giải quyết trực tiếp hai điểm yếu trên:
    \begin{enumerate}
        \item \textit{Cứu sống các truy vấn bị kênh ngữ nghĩa bỏ sót ngoài top 10}: Ở câu \texttt{e021} (hỏi về tiền tố thuộc tính nội bộ \texttt{\$ACTION\_} khi truyền đối tượng \texttt{FormData} của Next.js Server Action vào hàm \nolinkurl{Object.fromEntries}) và câu \texttt{e083} (hỏi về phương thức \nolinkurl{module.syncBuiltinESMExports()} trong Node.js), từ điển \texttt{simple} không khớp được token khiến nhánh Hybrid RRF đẩy mục đích ra ngoài top 10; nhưng kênh nguyên văn với \texttt{pattern = "\$ACTION"} và \texttt{pattern = "syncBuiltinESMExports"} đã bắt trúng đích trong giỏ \texttt{verbatim\_only}, nâng \texttt{Found} toàn cục từ $95,0\%$ lên \textbf{97,0\%}.
        \item \textit{Kéo ngược các phân đoạn dài bị Reranker hạ bậc về lại hạng 1}: Nhờ quy tắc ưu tiên giỏ giao \texttt{both} ($\mathcal{H}_{\text{sem}} \cap \mathcal{H}_{\text{verb}}$) lên trước \texttt{semantic\_only}, \textbf{10 câu hỏi} bị tầng Cross-Encoder làm tụt hạng ở cấu hình (4) gồm \texttt{e007} ($2 \to 1$), \texttt{e012} ($5 \to 1$), \texttt{e018} ($2 \to 1$), \texttt{e029} ($7 \to 1$), \texttt{e031} ($3 \to 1$), \texttt{e035} ($2 \to 1$), \texttt{e079} ($2 \to 1$), \texttt{e086} ($4 \to 1$), \texttt{e096} ($3 \to 1$) và \texttt{e098} ($4 \to 1$) đã được kéo thẳng về vị trí số 1 nhờ chứa đúng mẫu \texttt{pattern} (\texttt{listen(0)}, \texttt{lifo}, \texttt{cleanup}, \texttt{ticketKeys}).
    \end{enumerate}
    \item \textbf{Phân tích nguyên nhân gốc rễ của 3 ca bỏ sót (\texttt{e011}, \texttt{e070}, \texttt{e074}):} Cả 3 ca duy nhất có \texttt{found = False} trên toàn bộ 100 truy vấn đều không xuất phát từ thuật toán xếp hạng mà do \textbf{bộ lọc cứng \texttt{doc\_slugs} của tác tử Pha 2}: ở câu \texttt{e011} (hỏi về ngưỡng \texttt{cacheLife}), tác tử chỉ truyền 3 tài liệu \nolinkurl{08_caching}, \nolinkurl{incremental_static_regeneration_cache_components} và \nolinkurl{migrating_to_cache_components} trong khi đặc tả \texttt{cacheLife} nằm ở tài liệu \nolinkurl{09_revalidating}; tương tự, ở câu \texttt{e070} và \texttt{e074}, tác tử lọc vào các tài liệu lý thuyết \nolinkurl{prefetching} và \nolinkurl{rendering_philosophy} trong khi ví dụ thực hành nằm ở \nolinkurl{keeping_pages_static} và \nolinkurl{package_bundling}. Trong thực tế vận hành nhiều bước của \texttt{Coding Agent}, khi lượt gọi đầu tiên với \texttt{doc\_slugs} hẹp không trả về thông tin cần tìm, \texttt{Agent} sẽ bỏ bộ lọc \texttt{doc\_slugs} ở lượt gọi thứ hai để tìm trên toàn kho.
    \item \textbf{Hiệu năng vận hành thực tế của kiến trúc Warm MCP Server:} Khi tích hợp với các CLI Agent bên ngoài, việc khởi chạy tiến trình con \texttt{stdio} mới cho mỗi câu hỏi gây ra độ trễ khởi động lạnh lên tới $29.727\text{ ms}$ ở lệnh gọi \texttt{hybrid\_search} đầu tiên do phải nạp lại mô hình \texttt{Qwen3-Embedding-0.6B} và \texttt{CrossEncoderReranker} lên GPU. Bằng cách triển khai giao thức \textbf{Warm MCP Streamable HTTP} (\texttt{POST /api/mcp}) chia sẻ trực tiếp \texttt{RetrievalEngine} đã làm ấm sẵn trên FastAPI cùng cơ chế ưu tiên bộ nhớ đệm cục bộ (\texttt{local\_files\_only=True}), độ trễ lệnh gọi công cụ đầu tiên của \texttt{Agent} giảm xuống chỉ còn $695\text{--}768\text{ ms}$, đáp ứng trải nghiệm hỏi đáp thời gian thực.
\end{itemize}

\noindent\textbf{Tổng hợp trả lời các câu hỏi nghiên cứu (CH1--CH3):}
\begin{itemize}
    \item \textbf{Trả lời CH1:} Hai bộ \texttt{Parser} mẫu cho Next.js MDX và Node.js Markdown đã phân rã toàn bộ 167 tài liệu (162.757 dòng) thành 5.633 \texttt{chunk} trong $\approx 2,1$ giây, đạt \textbf{100\%} toàn vẹn ánh xạ dòng 1:1 (\texttt{line\_errors = 0}), \textbf{100\%} tiền tố ngữ cảnh hợp lệ (\texttt{ctx\_errors = 0}) và bảo toàn trọn vẹn 2.107 phân đoạn chứa code block (\texttt{KT01}, \texttt{KT02}).
    \item \textbf{Trả lời CH2:} Kết hợp Dense + Sparse qua RRF trên $\text{text}^{\text{ctx}}_i$ vượt trội có ý nghĩa thống kê so với Sparse FTS đơn lẻ ($p = 0,0156$), nâng \texttt{MRR@10} lên \textbf{79,0\%} và \texttt{Recall@5} lên \textbf{93,0\%}. Việc bổ sung kênh nguyên văn theo phân hoạch Venn cứu sống các truy vấn ký hiệu đặc thù (\texttt{e021}, \texttt{e083}), nâng tỷ lệ tìm thấy lên \textbf{97,0\%} (và \textbf{100\%} trên 97/97 câu có phạm vi lọc hợp lệ), đồng thời kéo 10 phân đoạn dài bị cửa sổ 256 token của Cross-Encoder hạ bậc về lại hạng 1.
    \item \textbf{Trả lời CH3:} Kiến trúc Warm Streamable HTTP MCP (\texttt{POST /api/mcp}) giảm độ trễ lệnh gọi công cụ đầu tiên từ $29.727\text{ ms}$ xuống $695\text{--}768\text{ ms}$ ($\approx 40\times$), trong khi công cụ \nolinkurl{submit_grounded_answer} đối chiếu trực tiếp từng bộ ba trích dẫn \texttt{(path, \allowbreak start\_line, \allowbreak end\_line)} trên PostgreSQL (\texttt{KT07}), chặn đứng mọi trích dẫn đường dẫn \texttt{ltree} không tồn tại trong kho ngữ liệu.
\end{itemize}
